\documentclass[11pt, a4paper]{article}\usepackage[]{graphicx}\usepackage[]{xcolor}
% maxwidth is the original width if it is less than linewidth
% otherwise use linewidth (to make sure the graphics do not exceed the margin)
\makeatletter
\def\maxwidth{ %
  \ifdim\Gin@nat@width>\linewidth
    \linewidth
  \else
    \Gin@nat@width
  \fi
}
\makeatother

\definecolor{fgcolor}{rgb}{0.345, 0.345, 0.345}
\newcommand{\hlnum}[1]{\textcolor[rgb]{0.686,0.059,0.569}{#1}}%
\newcommand{\hlsng}[1]{\textcolor[rgb]{0.192,0.494,0.8}{#1}}%
\newcommand{\hlcom}[1]{\textcolor[rgb]{0.678,0.584,0.686}{\textit{#1}}}%
\newcommand{\hlopt}[1]{\textcolor[rgb]{0,0,0}{#1}}%
\newcommand{\hldef}[1]{\textcolor[rgb]{0.345,0.345,0.345}{#1}}%
\newcommand{\hlkwa}[1]{\textcolor[rgb]{0.161,0.373,0.58}{\textbf{#1}}}%
\newcommand{\hlkwb}[1]{\textcolor[rgb]{0.69,0.353,0.396}{#1}}%
\newcommand{\hlkwc}[1]{\textcolor[rgb]{0.333,0.667,0.333}{#1}}%
\newcommand{\hlkwd}[1]{\textcolor[rgb]{0.737,0.353,0.396}{\textbf{#1}}}%
\let\hlipl\hlkwb

\usepackage{framed}
\makeatletter
\newenvironment{kframe}{%
 \def\at@end@of@kframe{}%
 \ifinner\ifhmode%
  \def\at@end@of@kframe{\end{minipage}}%
  \begin{minipage}{\columnwidth}%
 \fi\fi%
 \def\FrameCommand##1{\hskip\@totalleftmargin \hskip-\fboxsep
 \colorbox{shadecolor}{##1}\hskip-\fboxsep
     % There is no \\@totalrightmargin, so:
     \hskip-\linewidth \hskip-\@totalleftmargin \hskip\columnwidth}%
 \MakeFramed {\advance\hsize-\width
   \@totalleftmargin\z@ \linewidth\hsize
   \@setminipage}}%
 {\par\unskip\endMakeFramed%
 \at@end@of@kframe}
\makeatother

\definecolor{shadecolor}{rgb}{.97, .97, .97}
\definecolor{messagecolor}{rgb}{0, 0, 0}
\definecolor{warningcolor}{rgb}{1, 0, 1}
\definecolor{errorcolor}{rgb}{1, 0, 0}
\newenvironment{knitrout}{}{} % an empty environment to be redefined in TeX

\usepackage{alltt}

\usepackage[top = 0.6 in, bottom = 0.6 in, left = 0.8 in, right = 0.8 in]{geometry}
\usepackage{amsmath, amssymb, amsfonts}

\allowdisplaybreaks[4]

\usepackage{enumerate}
\usepackage{array}
\usepackage{multirow}
\usepackage{dingbat}
\usepackage{fontawesome5}
\usepackage{tasks}
\usepackage{bbding}
\usepackage{twemojis}
% how to use bull's eye ----- \scalebox{2.0}{\twemoji{bullseye}}
\usepackage{fontspec}


\title{MSMS 408 : Practical 04}
\author{Ananda Biswas \\[1em] Exam Roll No. : 24419STC053}
\date{\today}

\newfontface\myfont{Myfont1-Regular.ttf}[LetterSpace=0.05em]

\newfontface\cbfont{CaveatBrush-Regular.ttf}
% how to use --- \myfont --write text here--
\IfFileExists{upquote.sty}{\usepackage{upquote}}{}
\begin{document}

\maketitle


\section*{\faArrowAltCircleRight[regular] \textcolor{blue}{Question 1}}

\hspace{1cm} A hospital emergency helpline receives calls at an average rate of $3$ calls per hour. The calls occur randomly and independently over time. Let $X$ denote the number of calls received in a randomly selected hour. 

\begin{enumerate}[(i)]

\item Assuming a suitable probability model, identify the distribution of $X$ and write its probability mass function. 

\item Calculate the probability that exactly $4$ calls are received in an hour.

\item Determine the mean and variance of the distribution.

\item Generate simulated values of $X$ using the inverse transform method. 

\item Compare the simulated results with the theoretical probabilities and estimate the model parameter from the simulated data.

\end{enumerate}

\section*{\faArrowAltCircleRight[regular] \textcolor{blue}{Theory}}

\begin{itemize}
\item[\scalebox{2.0}{\twemoji{keycap: 1}}] \hspace{0.1cm} In the given set-up, $X$ is said to have a 
\end{itemize}

\section*{\faArrowAltCircleRight[regular] \textcolor{blue}{R Program}}

\begin{knitrout}
\definecolor{shadecolor}{rgb}{0.969, 0.969, 0.969}\color{fgcolor}\begin{kframe}
\begin{alltt}
\hlkwd{set.seed}\hldef{(}\hlnum{22}\hldef{)}
\end{alltt}
\end{kframe}
\end{knitrout}


\newpage

\section*{\faArrowAltCircleRight[regular] \textcolor{blue}{Question 2}}

\hspace{1cm} A marketing executive makes repeated phone calls to potential customers to sell a product. The probability that any individual call results in a successful sale is $0.3$, and each call is independent of the others. Let $X$ denote the number of calls required to achieve the first successful sale.

\begin{enumerate}[(i)]

\item Model $X$ using an appropriate probability distribution, write its probability mass function.

\item Compute the probability that the first success occurs on the $5$th call.

\item Determine the expected number of calls required for the first success. 

\item Using the inverse transform method, simulate multiple observations of $X$ and estimate the success probability from the simulated data.

\end{enumerate}

\section*{\faArrowAltCircleRight[regular] \textcolor{blue}{Theory}}



\section*{\faArrowAltCircleRight[regular] \textcolor{blue}{R Program}}





\end{document}
